% ============================================================
% NDH Manifold Definition
% Module: 01-ndh-manifold.tex
% Altitude: A1–A5 (Expressive → Rendering → Governance)
% Membrane: Manifold-Geometry • Constitutional Base-Space
% ============================================================

\section{The NDH Manifold}
\label{sec:ndh-manifold}

The NDH manifold is the geometric base-space upon which all constitutional
structures, operator adjacencies, and non-absorptive cognitive physics are
defined. It is a tenary manifold equipped with three adjacency categories
governing spectral, constitutional, and operator-level interactions.

% ------------------------------------------------------------
% 1. Base Space
% ------------------------------------------------------------

\subsection{Base Space}

Let $X$ be a second-countable, Hausdorff topological space. The NDH manifold
is defined as the triple

\[
\mathcal{M} = (X, \mathcal{A}_s, \mathcal{A}_c, \mathcal{A}_o)
\]

where:

\begin{itemize}
    \item $\mathcal{A}_s$ is the \emph{spectral adjacency category},
    \item $\mathcal{A}_c$ is the \emph{constitutional adjacency category},
    \item $\mathcal{A}_o$ is the \emph{operator adjacency category}.
\end{itemize}

Each adjacency category governs a distinct mode of geometric interaction.

% ------------------------------------------------------------
% 2. Spectral Adjacency
% ------------------------------------------------------------

\subsection{Spectral Adjacency}

Spectral adjacency encodes the manifold's modal structure. It is defined as
a small category

\[
\mathcal{A}_s = (\mathrm{Obj}_s, \mathrm{Hom}_s)
\]

whose objects correspond to spectral regions of $X$ and whose morphisms
represent permissible spectral transitions.

Spectral adjacency is curvature-sensitive and interacts with the Stability
Constitution.

% ------------------------------------------------------------
% 3. Constitutional Adjacency
% ------------------------------------------------------------

\subsection{Constitutional Adjacency}

Constitutional adjacency governs the behavior of constitutional membranes.
It is defined as

\[
\mathcal{A}_c = (\mathrm{Obj}_c, \mathrm{Hom}_c)
\]

where morphisms encode permissible constitutional transformations under the
Octo Constitution.

Constitutional adjacency is functorially constrained by the constitutional
stack $\mathfrak{C}$.

% ------------------------------------------------------------
% 4. Operator Adjacency
% ------------------------------------------------------------

\subsection{Operator Adjacency}

Operator adjacency governs the NDH operator-spine. It is defined as

\[
\mathcal{A}_o = (\mathrm{Obj}_o, \mathrm{Hom}_o)
\]

where morphisms correspond to operator-level transitions, including
constitutional operators, curvature operators, and identity-preserving
operators.

Operator adjacency interacts with the Operator Algebra module.

% ------------------------------------------------------------
% 5. Fiber Bundle Structure
% ------------------------------------------------------------

\subsection{Fiber Bundle Structure}

The NDH manifold carries a fiber bundle structure

\[
\pi : \mathcal{M} \to X
\]

where each fiber $\pi^{-1}(x)$ contains the adjacency categories associated
with the point $x \in X$.

This structure ensures that constitutional and operator adjacencies vary
smoothly across the manifold.

% ------------------------------------------------------------
% 6. Identity Preservation
% ------------------------------------------------------------

\subsection{Identity Preservation}

Identity is preserved across all adjacency morphisms. Formally,

\[
\forall f \in \mathrm{Hom}_s \cup \mathrm{Hom}_c \cup \mathrm{Hom}_o,\quad
\pi(f) = \mathrm{id}_X
\]

This constraint integrates with the Identity Constitution.

% ------------------------------------------------------------
% 7. Non-Absorptive Geometry
% ------------------------------------------------------------

\subsection{Non-Absorptive Geometry}

The NDH manifold is non-absorptive: no adjacency morphism may collapse,
inherit, or absorb external identity. This property is enforced by the
Quantum Sponge sheaf $\mathcal{Q}$ defined in Module 02.

% ------------------------------------------------------------
% End of Module
% ------------------------------------------------------------
